% ============================================================
%  TORPEDO AVANZADO HORIZONTAL – CÁLCULO I (INF-112) | UCM
%  Diseño de Ultra-Alta Densidad - 4 Páginas estrictas (LANDSCAPE)
% ============================================================
\documentclass[8pt,letterpaper,landscape]{extarticle}

% ── Paquetes Críticos de Configuración ───────────────────────
\usepackage[top=0.4cm,bottom=0.4cm,left=0.4cm,right=0.4cm,landscape]{geometry}
\usepackage{multicol}
\usepackage{amsmath,amssymb,mathtools}
\usepackage{titlesec}
\usepackage{xcolor}
\usepackage{booktabs}
\usepackage{array}
\usepackage{enumitem}
\usepackage{microtype}
\usepackage{parskip}
\usepackage[T1]{fontenc}
\usepackage{lmodern}

% ── Paleta de Colores Académica ──────────────────────────────
\definecolor{primary}{RGB}{15,45,100}
\definecolor{secondary}{RGB}{45,95,170}
\definecolor{accent}{RGB}{190,30,30}
\definecolor{lightbg}{RGB}{235,242,255}
\definecolor{darkgray}{RGB}{60,60,60}

% ── Ajustes de Espaciado y Layout Ultra-Compacto ─────────────
\setlength{\parskip}{0pt}
\setlength{\parindent}{0pt}
\setlength{\abovedisplayskip}{1pt}
\setlength{\belowdisplayskip}{1pt}
\setlength{\abovedisplayshortskip}{0pt}
\setlength{\belowdisplayshortskip}{0pt}
\setlength{\multicolsep}{2pt}
\setlength{\columnsep}{6pt}
\setlist{nosep,leftmargin=9pt,itemsep=0pt,topsep=0pt}
\pagestyle{empty}

% ── Macros de Bloques con Dimexpr ────────────────────────────
\newcommand{\bloque}[1]{%
  \vspace{2pt}%
  \noindent\colorbox{primary}{\parbox{\dimexpr\linewidth-2\fboxsep\relax}%
    {\color{white}\bfseries\fontsize{7.5pt}{8.5pt}\selectfont\space #1}}%
  \vspace{1pt}}

\newcommand{\subbloque}[1]{%
  \vspace{1pt}%
  \noindent\colorbox{lightbg}{\parbox{\dimexpr\linewidth-2\fboxsep\relax}%
    {\color{primary}\bfseries\fontsize{7pt}{8pt}\selectfont\space #1}}%
  \vspace{1pt}}

\newcommand{\cajadestacada}[1]{%
  \noindent\fcolorbox{accent}{white}{\parbox{\dimexpr\linewidth-2\fboxsep-2\fboxrule\relax}%
    {\fontsize{7pt}{8pt}\selectfont #1}}%
  \vspace{1pt}}

% ============================================================
\begin{document}
% ============================================================

% ============================================================
%  PÁGINA 1: LÓGICA Y DESIGUALDADES
% ============================================================
\noindent\colorbox{primary}{\parbox{\dimexpr\linewidth-2\fboxsep\relax}{%
  \centering\color{white}\bfseries\fontsize{8.5pt}{9.5pt}\selectfont
  INF-112 CÁLCULO I \quad|\quad PÁGINA 1: LÓGICA, VALOR ABSOLUTO Y SUPREMO \quad|\quad UCM}}
\vspace{2pt}

\begin{multicols}{4} % 4 columnas en horizontal aprovechan mejor el ancho

\bloque{1. Demostración por el Absurdo}
\subbloque{Esquema General para Irracionales ($\sqrt{p}$)}
\textbf{Objetivo:} Probar que si $p$ es primo, $\sqrt{p} \notin \mathbb{Q}$.
\begin{enumerate}
  \item \textbf{Hipótesis contraria:} Suponer que $\sqrt{p} \in \mathbb{Q}$. Existen $m, n \in \mathbb{Z}^+$ tales que $\sqrt{p} = \frac{m}{n}$ con $\gcd(m,n) = 1$ (fracción irreducible).
  \item \textbf{Álgebra:} Elevar al cuadrado: $p = \frac{m^2}{n^2} \implies p \cdot n^2 = m^2$.
  \item \textbf{Lema de Euclides:} Como $p \mid m^2$ y $p$ es primo $\implies p \mid m$. Por ende, existe $k \in \mathbb{Z}$ tal que $m = p \cdot k$.
  \item \textbf{Sustitución:} $p \cdot n^2 = (p \cdot k)^2 \implies p \cdot n^2 = p^2 \cdot k^2 \implies n^2 = p \cdot k^2$.
  \item \textbf{Segunda divisibilidad:} Se deduce que $p \mid n^2 \implies p \mid n$.
  \item \textbf{Contradicción:} $p \mid m$ y $p \mid n$, lo cual contradice la hipótesis de que $\gcd(m,n) = 1$. $\therefore \sqrt{p} \notin \mathbb{Q}$. $\blacksquare$
\end{enumerate}

\cajadestacada{\textbf{Ejemplo Examen:} Demostrar $\sqrt{23} \notin \mathbb{Q}$.\\
Supongamos $\sqrt{23}=\frac{m}{n}$ con $\gcd(m,n)=1$.\\
$23n^2 = m^2 \implies 23 \mid m \implies m = 23k$.\\
$23n^2 = 23^2 k^2 \implies n^2 = 23k^2 \implies 23 \mid n$.\\
Contradice $\gcd(m,n)=1$. Luego, $\sqrt{23}$ es irracional.}

\bloque{2. Valor Absoluto e Inecuaciones}
\subbloque{Definición y Propiedades}
$|A| = \begin{cases} A & \text{si } A \ge 0 \\ -A & \text{si } A < 0 \end{cases}$
\begin{itemize}
  \item $|A \cdot B| = |A| \cdot |B| \quad \bullet \quad \left|\frac{A}{B}\right| = \frac{|A|}{|B|}$
  \item $|A + B| \le |A| + |B|$ (Desigualdad Triangular).
  \item $|A| = \sqrt{A^2}$
\end{itemize}

\subbloque{Inecuaciones Fundamentales ($k > 0$)}
\begin{center}
\begin{tabular}{@{}ll@{}}
\toprule
\textbf{Inecuación} & \textbf{Forma Equivalente} \\ \midrule
$|A| \le k$ & $-k \le A \le k \iff A \ge -k \;\land\; A \le k$ \\
$|A| \ge k$ & $A \le -k \;\lor\; A \ge k$ \\ \bottomrule
\end{tabular}
\end{center}

\subbloque{Método de Puntos Críticos (Racionales)}
Para expresiones del tipo $\frac{P(x)}{Q(x)} \gtrless 0$:
\begin{enumerate}
  \item Factorizar por completo $P(x)$ y $Q(x)$ sobre $\mathbb{R}$.
  \item Hallar las raíces reales de $P(x)$ y $Q(x)$ (puntos críticos).
  \item Construir tabla de signos. \textcolor{accent}{\textbf{OJO:}} Los valores que anulan a $Q(x)$ llevan siempre intervalo abierto ($x \neq$ raíz de $Q$).
\end{enumerate}

\subbloque{Suma de Valores Absolutos}
Para resolver $|f(x)| + |g(x)| \le c$:
\begin{enumerate}
  \item Igualar a cero cada argumento: $f(x)=0 \land g(x)=0$.
  \item Definir los intervalos en la recta real.
  \item Redefinir la inecuación en cada intervalo removiendo las barras según su signo local. Resolver e intersectar con el intervalo analizado. $S_{\text{total}} = \bigcup S_i$.
\end{enumerate}

\bloque{3. Axioma del Supremo y Acotación}
\subbloque{Definiciones Globales ($A \subseteq \mathbb{R}, A \neq \emptyset$)}
\begin{itemize}
  \item \textbf{Cota Superior:} $M \in \mathbb{R}$ tal que $\forall a \in A, a \le M$.
  \item \textbf{Cota Inferior:} $m \in \mathbb{R}$ tal que $\forall a \in A, m \le a$.
  \item \textbf{Conjunto Acotado:} Posee cota superior e inferior.
\end{itemize}

\subbloque{Supremo ($\sup$) e Ínfimo ($\inf$)}
\begin{itemize}
  \item $\boldsymbol{s = \sup A}$: 1) $\forall a \in A, a \le s$. 2) $\forall \varepsilon > 0, \exists a_0 \in A : a_0 > s - \varepsilon$.
  \item $\boldsymbol{i = \inf A}$: 1) $\forall a \in A, i \le a$. 2) $\forall \varepsilon > 0, \exists a_0 \in A : a_0 < i + \varepsilon$.
\end{itemize}

\subbloque{Máximo ($\max$) y Mínimo ($\min$)}
\begin{itemize}
  \item $\max A = s \iff s = \sup A \;\land\; s \in A$.
  \item $\min A = i \iff i = \inf A \;\land\; i \in A$.
\end{itemize}

\cajadestacada{\textbf{Axioma del Supremo:} Todo subconjunto de $\mathbb{R}$ no vacío y acotado superiormente posee un supremo en $\mathbb{R}$.}

\end{multicols}
\newpage

% ============================================================
%  PÁGINA 2: GEOMETRÍA ANALÍTICA Y CÓNICAS
% ============================================================
\noindent\colorbox{secondary}{\parbox{\dimexpr\linewidth-2\fboxsep\relax}{%
  \centering\color{white}\bfseries\fontsize{8.5pt}{9.5pt}\selectfont
  INF-112 CÁLCULO I \quad|\quad PÁGINA 2: GEOMETRÍA ANALÍTICA Y CÓNICAS \quad|\quad UCM}}
\vspace{2pt}

\begin{multicols}{4}

\bloque{4. La Recta en el Plano}
\subbloque{Ecuaciones y Parámetros}
\begin{itemize}
  \item \textbf{Pendiente:} $m = \frac{y_2 - y_1}{x_2 - x_1}$
  \item \textbf{Punto-Pendiente:} $y - y_1 = m(x - x_1)$
  \item \textbf{General:} $Ax + By + C = 0 \implies m = -\frac{A}{B}$
\end{itemize}
\subbloque{Relaciones entre Rectas}
\begin{itemize}
  \item \textbf{Paralelas ($L_1 \parallel L_2$):} $m_1 = m_2$
  \item \textbf{Perpendiculares ($L_1 \perp L_2$):} $m_1 \cdot m_2 = -1 \implies m_2 = -\frac{1}{m_1}$
\end{itemize}

\bloque{5. Circunferencia}
\subbloque{Ecuación Canónica}
$(x - h)^2 + (y - k)^2 = r^2 \implies \boldsymbol{C(h,k)}, \boldsymbol{r=\text{radio}}.$
\subbloque{Ecuación General a Canónica}
$x^2 + y^2 + Dx + Ey + F = 0 \implies h = -\frac{D}{2}, \;\; k = -\frac{E}{2}$
\[ r = \frac{1}{2}\sqrt{D^2 + E^2 - 4F} \quad (\text{Requiere } D^2+E^2-4F > 0) \]

\subbloque{Intersección de dos Circunferencias}
Dadas $C_1: x^2+y^2+D_1x+E_1y+F_1 = 0$ y $C_2: x^2+y^2+D_2x+E_2y+F_2 = 0$:
\cajadestacada{\textbf{Método:} Restar $C_1 - C_2$. Los términos cuadráticos se eliminan generando el \textbf{Eje Radical}: $(D_1-D_2)x + (E_1-E_2)y + (F_1-F_2) = 0$.\\
Despejar una variable de esta recta, sustituirla en $C_1$ o $C_2$ y resolver la ecuación cuadrática resultante para hallar los puntos.}

\bloque{6. Completación de Cuadrados}
\cajadestacada{\[ ax^2 + bx = a\left(x + \frac{b}{2a}\right)^2 - \frac{b^2}{4a} \]}

\bloque{7. Secciones Cónicas}
\subbloque{1. Parábola}
\begin{center}
\begin{tabular}{@{}lll@{}}
\toprule
\textbf{Eje Focal} & \textbf{Ecuación Canónica} & \textbf{Dirección} \\ \midrule
Horiz. & $(y-k)^2 = 4p(x-h)$ & $p>0 \to \text{Der}, p<0 \to \text{Izq}$ \\
Vert. & $(x-h)^2 = 4p(y-k)$ & $p>0 \to \text{Arr}, p<0 \to \text{Aba}$ \\ \bottomrule
\end{tabular}
\end{center}
\begin{itemize}
  \item \textbf{Vértice:} $V(h,k)$. \textbf{Foco ($F$):} A distancia $|p|$ de $V$.
  \item \textbf{Directriz ($D$):} Recta perpendicular al eje focal a distancia $|p|$ de $V$ (lado opuesto al foco).
\end{itemize}

\subbloque{2. Elipse}
$\frac{(x-h)^2}{a^2} + \frac{(y-k)^2}{b^2} = 1 \quad \text{con } a > b > 0$. Relación: $\boldsymbol{a^2 = b^2 + c^2}$
\begin{itemize}
  \item \textbf{Centro:} $C(h,k)$. \textbf{Eje mayor:} Horizontal.
  \item \textbf{Vértices:} $V(h \pm a, k)$. \textbf{Focos:} $F(h \pm c, k)$.
  \item Si $b > a$, el eje mayor es vertical ($V(h, k \pm b)$ y $F(h, k \pm c)$ con $b^2 = a^2+c^2$).
\end{itemize}

\subbloque{3. Hipérbola}
\begin{itemize}
  \item \textbf{Eje Horiz.:} $\frac{(x-h)^2}{a^2} - \frac{(y-k)^2}{b^2} = 1$. Relación: $\boldsymbol{c^2 = a^2 + b^2}$.
  Focos: $F(h \pm c, k)$. Vértices: $V(h \pm a, k)$.
  \item \textbf{Eje Vert.:} $\frac{(y-k)^2}{a^2} - \frac{(x-h)^2}{b^2} = 1$. Focos: $F(h, k \pm c)$.
\end{itemize}
\subbloque{Asíntotas de la Hipérbola}
\cajadestacada{%
\textbf{Eje Horizontal:} $y - k = \pm \frac{b}{a}(x - h)$ \\
\textbf{Eje Vertical:} $y - k = \pm \frac{a}{b}(x - h)$%
}

\end{multicols}
\newpage

% ============================================================
%  PÁGINA 3: TRIGONOMETRÍA Y LÍMITES
% ============================================================
\noindent\colorbox{secondary}{\parbox{\dimexpr\linewidth-2\fboxsep\relax}{%
  \centering\color{white}\bfseries\fontsize{8.5pt}{9.5pt}\selectfont
  INF-112 CÁLCULO I \quad|\quad PÁGINA 3: TRIGONOMETRÍA, LÍMITES Y CONTINUIDAD \quad|\quad UCM}}
\vspace{2pt}

\begin{multicols}{4}

\bloque{8. Trigonometría Avanzada}
\subbloque{Identidades de Suma de Senos en Triángulos}
Si $\alpha + \beta + \gamma = \pi$ (ángulos interiores de un triángulo):
\cajadestacada{\[ \sin^2 \alpha + \sin^2 \beta + \sin^2 \gamma = 4\sin \alpha \sin \beta \sin \gamma \]}
\subbloque{Ley de Senos (Cerros y Funiculares)}
Uso clave en topografía y fuerzas coplanares:
\[ \frac{a}{\sin \alpha} = \frac{b}{\sin \beta} = \frac{c}{\sin \gamma} = 2R \]
Donde $R$ es el radio de la circunferencia circunscrita.

\bloque{9. Límites Notables}
Variables evaluadas cuando $x \to 0$:
\cajadestacada{%
\begin{itemize}
  \item $\displaystyle \lim_{x \to 0} \frac{\sin x}{x} = 1 \quad \bullet \quad \lim_{x \to 0} \frac{\tan x}{x} = 1$
  \item $\displaystyle \lim_{x \to 0} \frac{1 - \cos(kx)}{x^2} = \frac{k^2}{2}$
  \item $\displaystyle \lim_{x \to \infty} \left(1 + \frac{k}{x}\right)^x = e^k \quad \bullet \quad \lim_{x \to 0} \frac{e^x - 1}{x} = 1$
\end{itemize}%
}

\bloque{10. Técnicas Analíticas de Límites}
\subbloque{1. Racionalización (Conjugados)}
Para raíces algebraicas, multiplicar numerador y denominador por el conjugado: $(\sqrt{A} - \sqrt{B})(\sqrt{A} + \sqrt{B}) = A - B$.

\subbloque{2. División por la Mayor Potencia ($x \to \infty$)}
Dividir cada término por $x^n$, donde $n$ es el mayor grado del denominador. Si $\text{gr}(P) = \text{gr}(Q)$, el límite es el cociente de coeficientes principales.

\subbloque{3. Teorema del Sándwich (Acotamiento)}
Si $g(x) \le f(x) \le h(x)$ cerca de $c$, y además $\displaystyle \lim_{x \to c} g(x) = \lim_{x \to c} h(x) = L \implies \lim_{x \to c} f(x) = L$.
\textit{Uso común:} Funciones oscilantes acotadas (e.g., $-1 \le \sin(\frac{1}{x}) \le 1$).

\bloque{11. Continuidad Analítica}
\subbloque{Las 3 Condiciones Conceptuales}
$f(x)$ es continua en $x = c$ si y solo si:
\begin{enumerate}
  \item $f(c)$ existe.
  \item $\displaystyle \lim_{x \to c} f(x)$ existe $\iff \lim_{x \to c^-} f(x) = \lim_{x \to c^+} f(x)$.
  \item $\displaystyle \lim_{x \to c} f(x) = f(c)$.
\end{enumerate}

\subbloque{Parámetros $a$ y $b$ en Funciones por Partes}
Para garantizar continuidad en puntos de quiebre $x = c$:
\cajadestacada{%
\begin{enumerate}
  \item Calcular $\displaystyle \lim_{x \to c^-} f(x)$ (rama izquierda).
  \item Calcular $\displaystyle \lim_{x \to c^+} f(x)$ (rama derecha).
  \item Igualar límites laterales: $\lim_{x \to c^-} f(x) = \lim_{x \to c^+} f(x)$.
  \item Resolver el sistema de ecuaciones lineales resultante.
\end{enumerate}%
}

\end{multicols}
\newpage

% ============================================================
%  PÁGINA 4: DERIVADAS Y ASÍNTOTAS
% ============================================================
\noindent\colorbox{secondary}{\parbox{\dimexpr\linewidth-2\fboxsep\relax}{%
  \centering\color{white}\bfseries\fontsize{8.5pt}{9.5pt}\selectfont
  INF-112 CÁLCULO I \quad|\quad PÁGINA 4: DERIVADAS, ASÍNTOTAS Y OPTIMIZACIÓN \quad|\quad UCM}}
\vspace{2pt}

\begin{multicols}{4}

\bloque{12. Análisis de Asíntotas}
\subbloque{1. Asíntotas Verticales (AV)}
La recta $x = a$ es AV si $\displaystyle \lim_{x \to a^\pm} f(x) = \pm\infty$. Candidatos: Puntos que anulan al denominador pero no al numerador.

\subbloque{2. Asíntotas Horizontales (AH)}
La recta $y = L$ es AH si $\displaystyle \lim_{x \to \pm\infty} f(x) = L$ (con $L \in \mathbb{R}$).

\subbloque{3. Asíntotas Oblicuas (AO)}
Existen si $\text{gr}(\text{Num}) = \text{gr}(\text{Den}) + 1$. Tienen la forma $\boldsymbol{y = mx + b}$.
\cajadestacada{%
\[ m = \lim_{x \to \pm\infty} \frac{f(x)}{x} \quad (m \neq 0, m \in \mathbb{R}) \]
\[ b = \lim_{x \to \pm\infty} [f(x) - mx] \quad (b \in \mathbb{R}) \]%
}
\textit{Nota:} Si existe AH en una dirección, no hay AO en esa misma dirección.

\bloque{13. Cálculo de Derivadas}
\subbloque{Reglas Operacionales}
\begin{itemize}
  \item \textbf{Producto:} $(f \cdot g)' = f' \cdot g + f \cdot g'$
  \item \textbf{Cociente:} $\left(\frac{f}{g}\right)' = \frac{f' \cdot g - f \cdot g'}{g^2}$
  \item \textbf{Cadena:} $\frac{d}{dx}[f(g(x))] = f'(g(x)) \cdot g'(x)$
\end{itemize}

\subbloque{Derivadas Elementales}
\begin{tabular}{@{}ll|ll@{}}
\toprule
\textbf{$f(x)$} & \textbf{$f'(x)$} & \textbf{$f(x)$} & \textbf{$f'(x)$} \\ \midrule
$x^n$ & $n x^{n-1}$ & $e^x$ & $e^x$ \\
$\ln x$ & $\frac{1}{x}$ & $\sin x$ & $\cos x$ \\
$\cos x$ & $-\sin x$ & $\tan x$ & $\sec^2 x$ \\
$\arctan x$ & $\frac{1}{1+x^2}$ & $\arcsin x$ & $\frac{1}{\sqrt{1-x^2}}$ \\ \bottomrule
\end{tabular}

\bloque{14. Derivación Implícita}
Para curvas definidas por $F(x,y) = 0$:
\begin{enumerate}
  \item Derivar respecto a $x$ término a término.
  \item Recordar regla de la cadena para $y$: $\frac{d}{dx}(y^n) = n y^{n-1} \cdot y'$.
  \item Agrupar y despejar algebraicamente el término $y'$.
\end{enumerate}

\bloque{15. Recta Tangente y Normal}
La pendiente de la tangente es $\boldsymbol{m_t = f'(x_0)}$.
\cajadestacada{%
\textbf{Ecuación Recta Tangente:}
\[ y - y_0 = f'(x_0)(x - x_0) \]
\textbf{Ecuación Recta Normal (Perpendicular):}
\[ y - y_0 = -\frac{1}{f'(x_0)}(x - x_0) \]%
}

\bloque{16. Orden Superior y Criterios}
\begin{itemize}
  \item \textbf{Segunda Derivada ($f''(x)$):} Define concavidad.
  \item **Criterio de la 2da Derivada ($f'(x_0) = 0$):**
  1) Si $f''(x_0) < 0 \implies x_0$ es un **Máximo Local**.
  2) Si $f''(x_0) > 0 \implies x_0$ es un **Mínimo Local**.
  3) Si $f''(x_0) = 0 \implies$ No concluye (usar 1ra derivada).
\end{itemize}

\end{multicols}

% ── Pie de Página Oficial ────────────────────────────────────
\vfill
\noindent\rule{\linewidth}{0.4pt}
{\fontsize{6pt}{7pt}\selectfont\color{darkgray}
\textbf{INF-112 Cálculo I} – Torpedo de Ultra-Alta Densidad (Formato Horizontal). UCM. Diseñado para optimización de lectura bajo presión cronometrada. Impresión recomendada al 100\% de escala.}

% ============================================================
\end{document}
% ============================================================